\documentclass[11pt,a4paper]{article}
% Drake_preamble.tex - LaTeX Preamble Template
% 作者: 謝慕揚, MD, PhD, FESC
% 
% 使用說明:
% 1. 表格請使用 longtable，不要有垂直分隔線 |
% 2. 表格內換行使用 \newline，不要使用 \\
% 3. 藥物名稱、檢驗、檢查請維持英文
% 4. Reference 格式請使用 \href 加上驗證過的 DOI

% Including essential packages for Chinese typesetting
\usepackage{xeCJK}
\usepackage{fontspec}
% Configuring Chinese font with Noto Serif CJK TC, fallback to SimSun
\IfFontExistsTF{Noto Serif CJK TC}{
  \setCJKmainfont{Noto Serif CJK TC}
  \setCJKsansfont{Noto Serif CJK TC}
  \setCJKmonofont{Noto Serif CJK TC}
}{\setCJKmainfont{SimSun}
  \setCJKsansfont{SimSun}
  \setCJKmonofont{SimSun}
}

% Including other necessary packages
\usepackage{geometry}
\usepackage{titlesec}
\usepackage{enumitem}
\usepackage{xcolor}
\usepackage{colortbl}
\usepackage{tcolorbox}
\usepackage{booktabs}
\usepackage{longtable}
\usepackage{array}
\usepackage{graphicx}
\usepackage{tikz}
\usepackage{fancyhdr}
\usepackage{multicol}
\usepackage{datetime2}
\usepackage{hyperref}
\usepackage{mdframed}
\usepackage{amsmath}
\usepackage{amssymb}
\usepackage{multirow}

\hypersetup{
    colorlinks=true,
    linkcolor=emergency_blue,
    citecolor=emergency_blue,
    urlcolor=emergency_blue,
    pdfborder={0 0 0}
}

% Defining page geometry
\geometry{
  top=2.5cm,
  bottom=2.5cm,
  left=2cm,
  right=2cm,
  headheight=25pt
}

% Adjusting topmargin to compensate for increased headheight
\addtolength{\topmargin}{-10pt}

% Defining colors
\definecolor{emergency_red}{RGB}{186,24,27}
\definecolor{emergency_blue}{RGB}{0,114,188}
\definecolor{emergency_gray}{RGB}{120,120,120}
\definecolor{highlight_yellow}{RGB}{255,250,205}

% Setting title formats
\titleformat{\section}[block]{\Large\bfseries\color{emergency_red}}{\thesection}{10pt}{}
\titleformat{\subsection}[block]{\large\bfseries\color{emergency_blue}}{\thesubsection}{8pt}{}
\titleformat{\subsubsection}[block]{\normalsize\bfseries\color{emergency_gray}}{\thesubsubsection}{6pt}{}

% Defining custom environment for important notes
\newmdenv[
    linecolor=emergency_red,
    backgroundcolor=highlight_yellow,
    linewidth=2pt,
    topline=true,
    bottomline=true,
    leftline=true,
    rightline=true,
    innertopmargin=10pt,
    innerbottommargin=10pt,
    innerleftmargin=10pt,
    innerrightmargin=10pt
]{importantbox}

% Defining note command with tcolorbox
\newcommand{\note}[1]{
    \par
    \noindent
    \begin{tcolorbox}[
        colback=emergency_blue!5!white,
        colframe=emergency_blue,
        title=\textbf{重點提醒},
        fonttitle=\bfseries,
        boxrule=0.5pt,
        arc=3pt,
        left=6pt,
        right=6pt,
        top=6pt,
        bottom=6pt
    ]
    #1
    \end{tcolorbox}
}

% Key findings box
\newcommand{\keyfinding}[1]{
    \par
    \noindent
    \begin{tcolorbox}[
        colback=yellow!10!white,
        colframe=orange!75!black,
        title=\textbf{核心發現摘要},
        fonttitle=\bfseries,
        boxrule=0.5pt,
        arc=3pt
    ]
    #1
    \end{tcolorbox}
}

% Defining custom date format for ISO 8601
\DTMnewdatestyle{iso}{
  \renewcommand{\DTMdisplaydate}[4]{\number##1-\number##2-\number##3}
  \renewcommand{\DTMDisplaydate}[4]{\number##1-\number##2-\number##3}
}
\DTMsetdatestyle{iso}
\newcommand{\mydate}{\DTMtoday}


% Custom header
\pagestyle{fancy}
\fancyhf{}
\fancyhead[L]{\textcolor{emergency_red}{Critical Care 教學講義}}
\fancyhead[R]{\textcolor{emergency_blue}{ICUAW 懸吊式復健}}
\fancyfoot[C]{\textcolor{emergency_gray}{\thepage}}
\renewcommand{\headrulewidth}{1pt}
\renewcommand{\headrule}{\hbox to\headwidth{\color{emergency_red}\leaders\hrule height \headrulewidth\hfill}}

\begin{document}

%============================================================
% Title Page
%============================================================
\begin{center}
{\LARGE\bfseries\textcolor{emergency_red}{Critical Care}}\\[0.5cm]
{\Large\textbf{Efficacy of Suspension-Based Lower-Limb Rehabilitation Device in Enhancing Lower Limb Function Among Patients With ICU-Acquired Weakness}}\\[0.3cm]
{\large 懸吊式下肢復健裝置對 ICU 獲得性衰弱患者下肢功能之療效}\\[0.5cm]
{\normalsize \href{https://doi.org/10.1186/s13054-026-05840-1}{Crit Care (2026). DOI: 10.1186/s13054-026-05840-1}}\\[0.3cm]
{\normalsize 整理：謝慕揚 MD, PhD, FESC \quad 日期：\mydate}
\end{center}

\vspace{0.5cm}

%============================================================
\section{前言}
%============================================================

ICU 獲得性衰弱 (ICU-acquired weakness, ICUAW) 是重症患者常見且嚴重的併發症，發生率可達 48\%。ICUAW 導致脫離呼吸器困難、住院時間延長、功能障礙及生活品質下降。目前尚無實證臨床指引，主要策略為早期活動 (early mobilization) 及神經肌肉電刺激。本研究評估懸吊式下肢復健裝置 (suspension-based lower-limb rehabilitation device, SS) 結合傳統復健對 ICUAW 患者下肢功能的短期療效。

\keyfinding{
\begin{itemize}
    \item 本研究為 within-patient randomized controlled trial，以患者自身為對照
    \item 60 位 ICUAW 患者納入，54 位完成研究
    \item 主要結果：介入側肌肉厚度僅下降 0.2\%，對照側下降 6.3\%
    \item 次要結果：MRC score 及 AROM 均顯著改善
    \item 無介入相關不良事件
\end{itemize}
}

%============================================================
\section{研究設計}
%============================================================

\subsection{試驗設計}

\begin{itemize}
    \item \textbf{設計}：Prospective, within-patient randomized controlled trial
    \item \textbf{研究期間}：2023 年 4 月至 2025 年 5 月
    \item \textbf{研究地點}：上海長海醫院聯合 ICU 中心（共 116 床）
    \item \textbf{註冊}：ChiCTR2300070118
\end{itemize}

\subsection{納入與排除標準}

\begin{longtable}{p{3cm}p{11cm}}
\toprule
\textbf{標準} & \textbf{內容} \\
\midrule
\endfirsthead
\toprule
\textbf{標準} & \textbf{內容} \\
\midrule
\endhead
納入標準 & ICUAW 確診（MRC sum score $<$ 48）\newline 預期住院至少 2 週\newline 曾接受機械通氣至少 48 小時\newline 取得書面知情同意 \\
\midrule
排除標準 & 入 ICU 前或入住時心跳停止\newline Cardiac pacemaker 或 ICD\newline 急性腦損傷需深度鎮靜 $\geq$ 72 小時\newline Guillain-Barré syndrome 或重度失智\newline 未控制糖尿病、嚴重器官衰竭\newline 近期下肢 DVT 或 PE（治療 $<$ 48 小時）\newline 下肢骨折未癒合、嚴重骨質疏鬆\newline 現有神經肌肉疾病影響下肢功能 \\
\bottomrule
\end{longtable}

\subsection{介入方式}

\begin{longtable}{p{3cm}p{11cm}}
\toprule
\textbf{組別} & \textbf{介入內容} \\
\midrule
\endfirsthead
\toprule
\textbf{組別} & \textbf{介入內容} \\
\midrule
\endhead
對照側（Control） & 標準物理治療 40 分鐘/天\newline 包含：肌肉伸展、主被動關節活動度訓練、橫膈呼吸訓練、坐姿訓練\newline 每週 5 天，共 2 週 \\
\midrule
介入側（Intervention） & 標準物理治療 20 分鐘 + SS 訓練 20 分鐘/天\newline SS 裝置：SmartSling（上海卓道科技）\newline 提供 passive 及 active-assisted training\newline Hip/knee flexion-extension 及 adduction-abduction\newline 每週 5 天，共 2 週 \\
\bottomrule
\end{longtable}

\note{
SS 裝置為 flexible cable-driven 設計，可在床邊使用，提供多關節協調訓練。與固定式 cycling 裝置相比，SS 允許多平面運動及更大自由度，可能更有效促進神經肌肉活化。
}

%============================================================
\section{結果測量}
%============================================================

\subsection{主要結果}

\textbf{下肢肌肉厚度}（musculoskeletal ultrasound 測量）：
\begin{itemize}
    \item Rectus femoris (RF) + Vastus intermedius (VI) 合併厚度
    \item Vastus medialis (VM)
    \item Vastus lateralis (VL)
    \item Tibialis anterior (TA)
\end{itemize}

測量時機：Baseline、1 週、2 週

\subsection{次要結果}

\begin{itemize}
    \item \textbf{MRC score}：上下肢 6 組肌群肌力評估（0-60 分）
    \item \textbf{AROM}：Hip 及 knee 主動關節活動度
    \item \textbf{下肢周徑}：Thigh（髕骨上 10 cm）及 calf 周徑
\end{itemize}

%============================================================
\section{研究結果}
%============================================================

\subsection{受試者特徵}

\begin{itemize}
    \item 篩選 316 位，納入 60 位，54 位完成研究
    \item 平均年齡：68.7 $\pm$ 15.1 歲
    \item 男性：60\%
    \item APACHE II：17.2 $\pm$ 4.8
    \item SOFA：6.9 $\pm$ 2.3
    \item ICU 住院天數：55.9 $\pm$ 52.5 天
    \item 機械通氣天數：23.9 $\pm$ 19.0 天
    \item 主要診斷：感染 86.7\%（肺部 76.7\%）、ARDS 46.7\%
\end{itemize}

\subsection{主要結果：肌肉厚度}

\begin{longtable}{p{4cm}p{5cm}p{5cm}}
\toprule
\textbf{時間點} & \textbf{介入側變化} & \textbf{對照側變化} \\
\midrule
\endfirsthead
\toprule
\textbf{時間點} & \textbf{介入側變化} & \textbf{對照側變化} \\
\midrule
\endhead
1 週 & -0.2\% & -3.7\% \\
\midrule
2 週 & -0.1\% & -2.3\% \\
\midrule
\textbf{總計（2 週）} & \textbf{-0.2\%} & \textbf{-6.3\%} \\
\bottomrule
\end{longtable}

\textbf{統計分析}：
\begin{itemize}
    \item Group $\times$ Time interaction：F = 7.80, \textbf{p $<$ 0.001}
    \item Main effect of group：F = 10.87, \textbf{p = 0.001}
    \item 2 週時對照側顯著低於介入側（p $<$ 0.01）
    \item 介入側無顯著 within-group 變化（p = 1.00）
\end{itemize}

\subsection{次要結果}

\begin{longtable}{p{4cm}p{5cm}p{5cm}}
\toprule
\textbf{指標} & \textbf{Group $\times$ Time} & \textbf{組間差異} \\
\midrule
\endfirsthead
\toprule
\textbf{指標} & \textbf{Group $\times$ Time} & \textbf{組間差異} \\
\midrule
\endhead
MRC score & F = 4.46, p = 0.01 & 1 週 p = 0.05\newline 2 週 p = 0.003 \\
\midrule
AROM & F = 6.59, p = 0.001 & 1 週 p = 0.003\newline 2 週 p $<$ 0.001 \\
\midrule
下肢周徑 & F = 0.15, p = 0.86 & 無顯著差異 \\
\bottomrule
\end{longtable}

\subsection{亞組分析}

\begin{itemize}
    \item \textbf{年齡}：老年患者（$\geq$ 60 歲）獲益較大（p $<$ 0.001）
    \item \textbf{BMI}：正常 BMI 患者獲益較大（p = 0.004）
    \item \textbf{性別}：無顯著差異
\end{itemize}

\subsection{安全性}

\begin{itemize}
    \item 無介入相關不良事件
    \item 患者滿意度：91.7\% 滿意、3.3\% 尚可、5.0\% 未評估
\end{itemize}

%============================================================
\section{討論}
%============================================================

\subsection{主要發現}

SS 輔助復健結合傳統治療可有效減緩 ICUAW 患者的肌肉萎縮。雖然無法逆轉肌肉流失，但顯著減緩下降速度。MRC score 及 AROM 改善顯示 SS 可有效維持肌力及關節活動度。

\subsection{與 CYCLE Trial 比較}

\begin{longtable}{p{4cm}p{5cm}p{5cm}}
\toprule
\textbf{特點} & \textbf{本研究（SS）} & \textbf{CYCLE Trial} \\
\midrule
\endfirsthead
\toprule
\textbf{特點} & \textbf{本研究（SS）} & \textbf{CYCLE Trial} \\
\midrule
\endhead
裝置類型 & Flexible cable-driven & Fixed-track cycling \\
\midrule
運動平面 & 多平面（flexion、abduction） & 單一矢狀面 \\
\midrule
介入劑量 & 10 sessions / 2 週 & Median 3 sessions \\
\midrule
主要結果 & 正面（肌肉厚度保留） & 無顯著改善 \\
\bottomrule
\end{longtable}

\subsection{研究限制}

\begin{itemize}
    \item 單中心研究，需多中心驗證
    \item 短期追蹤（2 週），長期效果未知
    \item Within-patient 設計無法完全排除 cross-over effect
    \item 無法盲化治療師及患者
\end{itemize}

%============================================================
\section{Key Points}
%============================================================

\begin{enumerate}
    \item ICUAW 發生率可達 48\%，嚴重影響預後及生活品質
    \item SS 輔助復健結合傳統治療可減緩肌肉萎縮（介入側 -0.2\% vs 對照側 -6.3\%）
    \item MRC score 及 AROM 顯著改善，下肢周徑無顯著差異
    \item 老年患者及正常 BMI 患者獲益較大
    \item SS 裝置為 flexible cable-driven 設計，床邊使用便利，無相關不良事件
    \item 與 CYCLE trial（fixed-track cycling）相比，SS 的正面結果可能與多平面運動及較高介入劑量有關
    \item 此為初步結果，需大型多中心試驗確認
\end{enumerate}

%============================================================
\section{參考文獻}
%============================================================

\begin{enumerate}
    \item Xu L, Huang X, Wu H, et al. Efficacy of suspension-based lower-limb rehabilitation device in enhancing lower limb function among patients with ICU-acquired weakness: a self-controlled randomized clinical trial. \href{https://doi.org/10.1186/s13054-026-05840-1}{\textit{Crit Care}. 2026.}

    \item Fan E, Cheek F, Chlan L, et al. An official American Thoracic Society Clinical Practice guideline: the diagnosis of intensive care unit-acquired weakness in adults. \href{https://doi.org/10.1164/rccm.201411-2011ST}{\textit{Am J Respir Crit Care Med}. 2014;190(12):1437-46.}

    \item Hiser SL, Casey K, Nydahl P, et al. Intensive care unit acquired weakness and physical rehabilitation in the ICU. \href{https://doi.org/10.1136/bmj-2023-077292}{\textit{BMJ}. 2025;388:e077292.}

    \item Kho ME, Berney S, Pastva AM, et al. Early In-Bed Cycle Ergometry in Mechanically Ventilated Patients. \href{https://doi.org/10.1056/EVIDoa2400137}{\textit{NEJM Evidence}. 2024;3(7).}

    \item Schaller SJ, Scheffenbichler FT, Bein T, et al. Guideline on positioning and early mobilisation in the critically ill by an expert panel. \href{https://doi.org/10.1007/s00134-024-07532-2}{\textit{Intensive Care Med}. 2024;50(8):1211-1227.}
\end{enumerate}

\end{document}
